\section{Why the Responses Vary: Conditions and Mechanism}
\label{sec:why}

The observations of \S\ref{sec:observed} differ by method, capability, student and distribution. This
section asks what shapes that variation, and sorts the answers into causes that were verified,
explanations that were excluded, and factors that remain unidentified.

\subsection{Function shape}
\label{sec:why_shape}

Under pruning the pre-cliff response is a power law in the deleted capability mass once the measured
first-order alignment term is removed, with in-sample $R^2$ from 0.82 to 0.996 over 24
capability-by-model fits and a shared exponent that transfers across families (leave-one-family-out
MAE 2.17 nats against 10.3 for raw sparsity; Appendix~\ref{sec:law}); the cliff itself is a geometric
threshold crossing (Appendix~\ref{sec:app_mech}). Under quantization the response is a threshold in
bit-width whose position depends on the group size, which is why an interpolation over the measured
grid predicts and a smooth surface does not.

Under distillation the shape question was sharper. A four-coefficient joint form with a logarithmic
reuse term, $\delta=(a+a'z)\log(1+T/T_{\mathrm{ref}})+(b+b'z)\log(1+E)$, carries a falsifiable
implication: the pool-change effect divided by its log bracket must be flat in budget. It is not; the
ratio drifts by a factor of twenty across checkpoints for code on the 1B student, and the drift
survives at budget-matching tolerances of 1\% and 0.3\%. Opening the reuse term to a free exponent
$h_p(E)=((1+E)^p-1)/p$, profiled inside training folds, returns $p$ near or above one for every
capability (code 0.69, math 1.12, QA 1.44) with wide fold ranges and occasional boundary hits, and buys
no held-out accuracy: against six baselines on four grouped hold-outs, neither the log form, the
curved form nor a single interaction term beat the strongest development-stage baseline on either
pre-declared target (Appendix~\ref{sec:app_forms}).

Whether the failure lay in additivity itself could not be read from that data, because a fixed pool
and a doubling budget grid never place two trajectories on the same budget, so no complete
budget-by-reuse rectangle existed. Four trajectories were designed by replaying the trainer's update
boundaries to land on one, and did so to the token. The second difference across the rectangle, which
is exactly zero for any additive structure whatever its curvature, was $-0.24$ nats on the 1B student
and $-0.21$ on the 4B for the QA probe, inside the pre-registered band of $\pm 0.32$ and between the
additive prediction of zero and the frozen interaction candidate's $-0.49$
(Fig.~\ref{fig:explanation}b). The pre-fixed rule did not reject additivity; it also does not
establish it, and the 4B student's statistic on math and code fell outside its band at readouts the
design had registered as underpowered.

\subsection{Source state}
\label{sec:why_source}

Size and training stage of the source enter the pruning amplitude and carry math and code to unseen
densities of seen states; they over-predict QA on every new checkpoint, whose density response changes
sign and is not proportional to any shared curve ($R^2$ 0.49 against 0.94 and 0.97 for math and
code). Richer descriptors did not repair transfer: dense pre-compression statistics improve a
parametric family in four of nine arm-and-capability pairs and lose to the source-free median
configuration curve in nine of nine, and a heavily shrunk source correction added to that frozen curve
gains on none of the nine primary pairs (Appendix~\ref{sec:app_inputs}). Within the inputs and models
tested, the state helps where the response shape is shared across states and not where it is not.

\subsection{Training protocol and data reuse}
\label{sec:why_protocol}

Distillation damage on math and code is monotone in the learning rate: at $5\times10^{-5}$ the 1B
student is essentially unharmed ($+0.003$ nats) while still gaining 0.79 on QA, and at
$1\times10^{-4}$ the damage is an order of magnitude larger (0.051 on math, 0.052 on code); the size
ordering survives at every rate. Reported
damage magnitudes are therefore properties of the protocol as much as of distillation, and every
claim here names its rate. At a fixed rate the independent-data axis dominates: the reuse cost is
monotone in the reuse ratio $E=T/D_U$, grows with student size, and is paid mostly by QA
(\S\ref{sec:observed}).

\subsection{Evaluation distribution}
\label{sec:why_dist}

The cost of over-reuse transfers and the gain does not. All three QA distributions worsen at high
reuse and improve as reuse falls, for every student; only the 2Wiki sample crosses into a gain, while
MuSiQue flattens near zero and TriviaQA stays positive at every reuse level. A requirement phrased as
the independent data needed to keep a conditional loss increase below a threshold is therefore
supportable in form across distributions; a claim that distillation improves question answering is a
2Wiki statement.

\subsection{Mechanism}
\label{sec:why_mech}

For pruning and quantization the loss change under a logit displacement $r$ obeys an exact identity,
$\Delta L=(\mathbb{E}_p[r]-r_y)+\tfrac12\mathrm{Var}_p(r)+O(\lVert r\rVert^3)$, and its coefficient-free
second-order truncation reproduces measured damage with median relative error under 5\% up to a
boundary set by damage magnitude, not by family or knob: grouped four-bit quantization with fine
groups at 0.06 nats and 3.3\% error sits on the same curve as per-channel four-bit at 0.24 nats and
23.8\% (Fig.~\ref{fig:explanation}c, Appendix~\ref{sec:app_disp}). A shift-invariant projection of
the displacement onto the logits shows the pure-shrinkage account failing on magnitude, with about
three percent of the response in the shrinkage direction. What this account does not do is predict:
the variance of the orthogonal residual requires the compressed model, so it explains the damage
after the fact and is not an admitted input.

\subsection{What is settled, excluded, and open}
\label{sec:why_summary}

\emph{Verified causes.} The learning rate sets distillation damage magnitude. The reuse ratio sets
the QA cost, monotonically and with student size. The response shape differs across capabilities
under pruning, which is where capability conditioning earns its 0.06 nats macro (0.17 on QA).
Compression damage is second-order in logit displacement inside a magnitude boundary.

\emph{Excluded explanations.} A logarithmic reuse term with budget-independent coefficients. Pure
logit shrinkage as the damage mechanism. Source descriptors as a general route to cross-source
transfer within the inputs tested.

\emph{Unidentified.} Whether budget and reuse interact in the distillation response: the QA test
did not reject additivity and did not establish it, and neither the additive nor the frozen
interacting structure predicted the measured second difference within the registered noise. Why the
4B student's second differences on math and code exceed both structures' predictions while the 1B's
do not. The form of the reuse term beyond ``closer to linear than logarithmic''. What predicts the
data-change intervention for any capability, since nothing tested does.
